\fancyhead[L]{\small BESSEL ON THE COMET OF 1807.}
\hypertarget{native-535lower}{}\label{native-535lower}
\begin{center}\small Göttingische gelehrte Anzeigen. No. 26. pp. 250..256. February 16, 1811.\\[3pt]\rule{.72\linewidth}{.3pt}\end{center}

\emph{Investigations into the Apparent and True Orbit of the Great Comet That Appeared in \textup{1807}, by} F. W. \textsc{Bessel}, professor of astronomy in Königsberg. Published by \textsc{Friedr. Nicolovius}. Königsberg. 82 pages in quarto.

The comet of 1807, brighter than any that had been visible for many years, was, as is well known, regarded as one of the most remarkable celestial phenomena. It was of still greater interest to astronomers because it remained visible for so long. The observers, as if vying with one another, took advantage of this to determine its apparent motion as fully and precisely as possible. Several astronomers also undertook to calculate its parabolic orbit. No one, however, has devoted as much care to investigating the theory of this comet as the author of the present work, whose distinguished contributions on several earlier comets, especially that of 1769, are already well known. Much of the author's research on the comet of 1807 had already been made known through the Monatliche Crrespondenz and the Astronomische Jahrbuch. The purpose of the present work is to give a connected account of all this together with the later investigations by which he put the finishing touches to the theory of this comet. The work provides fresh evidence both of the present state of the art of observation and of the power of calculation when guided by so skilful a hand.

The \emph{first section} deals with the comet's apparent path across the sky. Mr. \textsc{Bessel} had no intention of collecting all the observations of the comet. He meant to bring together only as many carefully made observations as were needed to give an adequate account of the comet's apparent motion throughout its visibility from September 22, 1807, to March 27, 1808. He therefore begins by presenting his own observations made with the ring micrometer at Lilienthal. They deserve particular confidence given the author's conscientious care and his exceptional skill in this kind of work. Many readers will welcome the various precautions and remarks on this kind of observation that he gives here. \textsc{Bessel} determined the diameter of the field of view using pairs of stars whose declinations are known exactly and differ by nearly that diameter. This method is much to be recommended, and one can then express the diameter of the field by a simple series, as Bessel does. Whenever the diameter exceeds the difference between the declinations of the stars by more than one minute of arc, however, the reviewer prefers to use the formulas \[
\frac{a'-a}{2d}=\mathrm{tang}\,\varphi,\qquad \frac{a'+a}{2d}=\mathrm{tang}\,\psi,\qquad \frac{d}{\cos\varphi\cos\psi}=D
\] where $a$ and $a'$ denote the two chords, $d$ the difference between the declinations of the stars, and $D$ the diameter of the field of view. The observations made by \textsc{Bessel} himself extend from October 4, 1807, to February 24, \SourceLeaf{536}1808. In presenting them he has followed the highly commendable practice of including enough detail to allow them to be reduced again when the comparison stars, taken for the most part from the Histoire céleste, are known more accurately. Next come the observations of Dr. \textsc{Olbers} at Bremen, extending from October 8, 1807, to February 14, 1808, followed by those of \textsc{Thulis} at Marseille from September 22 to October 2 and those of \textsc{Wisniewsky} at Petersburg from March 18 to 27. The author reduced both of these latter sets afresh. Finally, to fill the gap that still remained in February 1808, there are the observations by \textsc{Oriani} at Milan from February 13 to 28.

The \emph{second section} is devoted to determining the comet's true orbit and is itself divided into three subdivisions. The \emph{first subdivision} gives parabolic and purely elliptical elements for the comet. The author calculated the first parabolic elements as early as October 23. Further observations already allowed him to improve them on November 6, and in March 1808 he calculated a third system of parabolic elements that fitted as closely as possible the earliest and latest observations then known to him and the right ascensions of the intervening observations. In the intervening declinations, however, a difference of one minute of arc still remained. It could have been spread to a considerable extent over the other five quantities, but would still have been too large to avoid concluding that there was a perceptible departure from the parabola. \textsc{Bessel} accordingly abandoned the parabolic hypothesis and, independently of that hypothesis, determined a fourth system of elements, which give an ellipse with an orbital period of 1953 years. When compared with all the Lilienthal and Bremen observations, these elements nowhere showed any differences except ones that varied so irregularly that they could only be attributed to errors in the individual observations. Soon afterward, however, the first Marseille observations showed that those elements gave right ascensions that were slightly too large. This led him to calculate a fifth system of elements, which reduced the orbital period to 1483 years. These elements represented all 91 observations so well that nowhere could a definite difference still be stated with confidence. In this way the orbit was determined as accurately as was possible using purely elliptical elements.

It was to be expected, however, that taking into account the perturbations that the comet must have undergone from the planets of our solar system while it was visible would modify these results all the more substantially because the orbital period depends so strongly on the small difference between the eccentricity and unity. A change in that difference, even one small in itself, can alter the orbital period enormously. It was therefore an interesting undertaking to allow properly for these perturbations, since otherwise it was quite impossible to judge within what limits the orbital period could be confined. The next two subdivisions serve this purpose. In the \emph{second subdivision}, the author develops the general method for calculating the comet's perturbations. As is well known, the most convenient way to do this is to treat the comet's orbit as a varying ellipse, to compute numerically the daily changes in each of the quantities that determine it under the action of the perturbing forces, and to determine their accumulated effect by mechanical quadratures. Several geometers have already derived the formulas for the instantaneous change in the elements of an elliptical orbit under perturbing forces. The reader will nevertheless take all the more pleasure in the new derivation that \textsc{Bessel} gives in this subdivision with great analytical elegance, since some modification was necessary for an ellipse of such high eccentricity. Specifically, the time of passage through perihelion was chosen in place of the epoch of the mean longitude or anomaly. The \SourceLeaf{537}expression readily obtained for the instantaneous perturbation of that time was ill suited to computation owing to the circumstance just mentioned and first required a special transformation.

In the \emph{third subdivision}, the author finally applies this method to the comet of 1807, first calculating the sum of the perturbing forces for seven equally spaced times from September 22, 1807, to March 20, 1808. He does not say explicitly whether he took all the planets into account or how many he considered. The reviewer would have liked to find the contribution of each planet stated separately, since a view of their relative contributions would not have been without interest. He next calculates the amount of the daily change in each of the six elements at those seven times. Taking his last system of elements to hold for September 22, 1807, he first mechanically integrates their instantaneous changes and then calculates their values at six different times from September 28, 1807, to March 23, 1808. For each of these times he had also determined the geocentric position with the greatest care from all the nearby observed positions. Comparing these positions with the calculations from those varying elements, and using the numerically derived differential relations between the individual elements and the calculated geocentric longitudes and latitudes, gave him 12 linear equations in the corrections to the elements still to be determined. From these, the values of the corrections were derived by the method of least squares, with due regard for the relative weight of the individual positions. This finally gave the definitive elements for September 22, of which we give here only the semimajor axis $= 143{,}195$ and the orbital period $= 1713\tfrac{1}{2}$ years. This orbital period, however, is continually altered by the ongoing perturbations. Even if this result were otherwise completely reliable, the time of the actual return could still be given only after an endlessly lengthy calculation. From a rough estimate of the perturbations due to Jupiter in the years 1808..1815, beside which the perturbations due to the other planets become imperceptible, \textsc{Bessel} already finds a shortening of the orbital period by 170 years.

In closing, the author also investigates the degree of accuracy that may be assigned to these results. Suppose that the errors in longitude and latitude are $5''$ at the first three normal places on September 28, October 22, and November 11, twice as large at the places on December 8 and February 21, and four times as large at the place on March 23. In their most unfavourable combination the orbital period could be increased to 2157 years or decreased to 1404 years. Such a combination is possible, though highly improbable. Given how carefully \textsc{Bessel} used so many observations to determine the normal places, there is also hardly any reason to fear that errors larger than those just assumed should have remained in them. There is therefore good reason to hope that the orbital period determined above is accurate to within a hundred years and that the comet will return around the thirty-fourth century of our era. Our distant descendants will then remember \textsc{Bessel}'s work with gratitude.

These last considerations also readily lead to the complete conviction that the comet \emph{certainly} describes an ellipse, and neither a parabola nor a hyperbola. To regard a non-elliptical orbit as possible, one would have to allow errors in the normal places \emph{at least} seven times as large, which can be declared outright impossible without hesitation. Thus, given the present perfection of \emph{practical astronomy}, the comet of 1807 is one of the few of which we can assert with certainty that they will certainly return one day, if no forces other than those we know act on them.
